\chapter{Casos de uso detallados}
\label{app:casos-uso}

Este apéndice recolle a especificación detallada dos casos de uso
representados na Figura~\ref{fig:diagrama-casos-uso} do
Capítulo~\ref{chap:Analise}. Para cada caso indícanse o obxectivo,
os actores participantes, o disparador, as precondicións,
as postcondicións e os fluxos que pode seguir a interacción.


% ============================================================
% CU-01
% ============================================================

\section{CU-01: Configurar a aplicación por primeira vez}
\label{cu:configuracion-inicial}

\begin{description}

\item[Obxectivo:] Preparar unha instalación nova da aplicación para que poida
empregarse como paciente ou como supervisor.

\item[Actor principal:] Paciente ou supervisor.

\item[Actor secundario:] Firebase Authentication.

\item[Disparador:] A persoa usuaria abre a aplicación por primeira vez ou nunha
instalación na que a configuración aínda non foi finalizada.

\item[Precondicións:] A aplicación está instalada e existe conexión á rede para
realizar a autenticación anónima.

\item[Fluxo principal:] \hfill
\begin{enumerate}
    \item A aplicación comproba que a configuración inicial aínda non foi
    finalizada e mostra a pantalla de benvida.

    \item A persoa usuaria escolle se vai xestionar o seu propio tratamento,
    actuando como paciente, ou o tratamento doutra persoa, actuando como
    supervisor.

    \item A aplicación realiza unha autenticación anónima mediante Firebase
    Authentication.

    \item Se a autenticación se completa correctamente, garda localmente o rol
    seleccionado.

    \item Se se escolle o rol de paciente, a aplicación mostra o proceso de
    vinculación cun supervisor. O paciente pode completar o CU-02 ou saltar
    este paso.

    \item O paciente introduce, se o desexa, un nome e establece a hora habitual
    da toma. A hora dispón inicialmente dun valor predefinido e pode ser
    modificada.

    \item A aplicación garda estes datos e marca a configuración do paciente
    como finalizada.

    \item Se se escolle o rol de supervisor, a aplicación abre directamente
    o fluxo de vinculación do CU-02. Neste caso debe completarse unha primeira
    vinculación antes de finalizar a configuración.

    \item Unha vez finalizada a configuración, ábrese a pantalla principal
    correspondente ao rol seleccionado.
\end{enumerate}

\item[Fluxos alternativos:] Se falla a autenticación anónima, móstrase unha
mensaxe de erro e non se continúa coa configuración. O paciente pode saltar a
vinculación inicial e utilizar a aplicación de maneira autónoma. O supervisor
non pode completar a súa configuración inicial sen vincular polo menos un
paciente.

\item[Postcondicións:] O rol queda almacenado no dispositivo e a instalación
queda preparada para acceder á pantalla principal correspondente. No caso do
paciente tamén quedan gardados a hora habitual da toma e, se se indicou, o nome.

\item[Requisitos relacionados:] RF-01, RF-02 e RF-03.

\end{description}


% ============================================================
% CU-02
% ============================================================

\section{CU-02: Vincular paciente e supervisor}
\label{cu:vincular-paciente-supervisor}

\begin{description}

\item[Obxectivo:] Establecer unha relación segura entre o dispositivo dun
paciente e o dunha persoa supervisora para permitir o intercambio de información
relacionada co seguimento do tratamento.

\item[Actores principais:] Paciente e supervisor.

\item[Actor secundario:] Firebase Cloud Messaging.

\item[Disparador:] O paciente decide mostrar o seu código QR e o supervisor
inicia a vinculación dun paciente.

\item[Precondicións:] Ambos dispositivos dispoñen dunha sesión anónima e dos
roles correspondentes. É necesaria conexión á rede e o supervisor debe poder
utilizar a cámara do dispositivo.

\item[Fluxo principal:] \hfill
\begin{enumerate}
    \item O dispositivo do paciente obtén o seu identificador e o token necesario
    para recibir mensaxes.

    \item A aplicación xera os datos de vinculación e crea un código QR que os
    representa.

    \item O paciente mostra o código QR ao supervisor.

    \item O supervisor abre o lector da aplicación e escanea o código QR.

    \item A aplicación comproba que o servizo necesario para realizar a
    comunicación está dispoñible e valida os datos recibidos.

    \item O dispositivo do supervisor obtén os seus datos de identificación e
    xera a información criptográfica necesaria para establecer a vinculación
    definitiva.

    \item O supervisor envía ao dispositivo do paciente unha mensaxe de
    vinculación protexida mediante a información contida no QR.

    \item O paciente recibe a mensaxe, comproba a súa validez e completa a
    vinculación segura.

    \item Ambos dispositivos gardan localmente a información necesaria para
    identificar a relación establecida.

    \item O paciente envía ao novo supervisor un resumo actualizado do seu estado
    para inicializar o seguimento.
\end{enumerate}

\item[Fluxos alternativos:] Se non se pode obter algún dos identificadores ou
tokens necesarios, o código QR non é válido, a conexión non está dispoñible ou
non se pode entregar a mensaxe inicial, a vinculación non se completa. O
paciente pode xerar novamente o código e o supervisor pode repetir o escaneo.
Durante a configuración inicial, o paciente tamén pode continuar sen establecer
unha vinculación.

\item[Postcondicións:] Os dous dispositivos conservan unha relación segura que
permite intercambiar información. Un paciente pode manter varias persoas
supervisoras vinculadas e un supervisor pode manter varios pacientes.

\item[Requisitos relacionados:] RF-04, RF-05, RF-19 e RF-24.

\end{description}


% ============================================================
% CU-03
% ============================================================

\section{CU-03: Configurar preferencias}
\label{cu:configurar-preferencias}

\begin{description}

\item[Obxectivo:] Permitir ao paciente modificar as preferencias empregadas
durante o seguimento do tratamento.

\item[Actor principal:] Paciente.

\item[Actor secundario:] Firebase Cloud Messaging, cando existen supervisores
vinculados.

\item[Disparador:] O paciente accede á pantalla de axustes da aplicación.

\item[Precondicións:] A aplicación está configurada co rol de paciente.

\item[Fluxo principal:] \hfill
\begin{enumerate}
    \item A aplicación carga o nome gardado, a hora habitual da toma e o estado
    do modo sinxelo.

    \item O paciente pode modificar o nome, que é opcional.

    \item O paciente pode seleccionar unha nova hora habitual para a toma.

    \item O paciente pode activar ou desactivar o modo sinxelo.

    \item Se solicita desactivar o modo sinxelo, a aplicación pide confirmación
    antes de volver mostrar todas as opcións da interface.

    \item O paciente confirma o gardado dos cambios.

    \item A aplicación garda o nome, a hora e o estado do modo sinxelo. Se o nome
    queda baleiro, non se conserva ningún nome asociado.

    \item Os recordatorios locais das tomas son reprogramados segundo a nova
    configuración.

    \item Se existen supervisores vinculados, envíaselles o estado actualizado
    do paciente.

    \item Ao regresar á pantalla principal, a aplicación aplica as novas
    preferencias.
\end{enumerate}

\item[Fluxos alternativos:] Se o paciente cancela a selección dunha nova hora,
mantense a anterior. Se cancela a desactivación do modo sinxelo, este permanece
activo. Se sae da pantalla sen gardar, os cambios non se aplican. Se non se pode
comunicar a actualización aos supervisores, as preferencias gardadas localmente
continúan sendo válidas.

\item[Postcondicións:] As preferencias seleccionadas quedan almacenadas no
dispositivo e os recordatorios quedan adaptados á hora configurada.

\item[Requisitos relacionados:] RF-06, RF-18 e RF-19.

\end{description}


% ============================================================
% CU-04
% ============================================================

\section{CU-04: Dixitalizar unha folla de tratamento}
\label{cu:dixitalizar-folla}

\begin{description}

\item[Obxectivo:] Obter automaticamente os datos estruturados dunha folla de
tratamento a partir dun documento proporcionado pola persoa usuaria.

\item[Actor principal:] Paciente ou supervisor.

\item[Actor secundario:] Servizo de procesado automático de follas de
tratamento.

\item[Disparador:] O paciente escolle engadir unha nova folla ou o supervisor
inicia o procesamento dunha folla para un paciente vinculado.

\item[Precondicións:] Existe conexión á rede e o documento está dispoñible para
ser fotografado ou seleccionado. Se o actor é un supervisor, debe existir un
paciente vinculado seleccionado.

\item[Fluxo principal:] \hfill
\begin{enumerate}
    \item A aplicación ofrece tres posibilidades: realizar unha fotografía,
    seleccionar unha imaxe da galería ou escoller un ficheiro PDF.

    \item A persoa usuaria selecciona unha das opcións e proporciona o documento.

    \item A aplicación envía o ficheiro ao servidor para o seu procesamento.

    \item O servidor comproba que o documento pode ser tratado e realiza, cando
    corresponde, o preprocesamento necesario.

    \item O documento é enviado ao servizo de procesado automático de follas de
    tratamento.

    \item O sistema comproba que o resultado obtido é suficientemente fiable e
    que contén a información necesaria para reconstruír a pauta.

    \item Os datos extraídos transfórmanse á estrutura empregada pola aplicación.

    \item O cliente recibe a cabeceira da folla, a pauta de doses e o histórico
    dispoñible.

    \item A aplicación abre o CU-05 para que a persoa usuaria revise o resultado
    antes de incorporalo ao tratamento.
\end{enumerate}

\item[Fluxos alternativos:] Se a persoa usuaria cancela a captura ou selección,
non se realiza ningún cambio. Se o formato non é admitido, o documento non pode
ser procesado, non se detecta correctamente a información necesaria ou se
produce un erro de conexión, móstrase unha mensaxe de erro e pode repetirse a
operación con outro documento.

\item[Postcondicións:] Os datos extraídos quedan dispoñibles temporalmente para
a súa revisión. A pauta previamente gardada non se modifica.

\item[Requisitos relacionados:] RF-08, RF-09, RF-10 e RF-21.

\end{description}


% ============================================================
% CU-05
% ============================================================

\section{CU-05: Revisar, corrixir e confirmar unha pauta}
\label{cu:revisar-confirmar-pauta}

\begin{description}

\item[Obxectivo:] Permitir comprobar e corrixir os datos extraídos antes de
incorporar unha nova pauta ao dispositivo do paciente.

\item[Actor principal:] Paciente ou supervisor.

\item[Actor secundario:] Firebase Cloud Messaging, cando existen dispositivos
vinculados implicados na operación.

\item[Disparador:] Finaliza correctamente o procesamento dunha folla no CU-04.

\item[Precondicións:] Existe un resultado da extracción pendente de revisión.
Se o actor é un supervisor, a folla está asociada a un paciente vinculado.

\item[Fluxo principal:] \hfill
\begin{enumerate}
    \item A aplicación ordena e mostra a información extraída da folla de
    tratamento.

    \item A persoa usuaria revisa as datas, as doses, os días de control e a
    seguinte visita.

    \item Se detecta algún erro, accede ao modo de corrección.

    \item A persoa usuaria pode modificar a data ou a dose dunha fila, indicar
    se corresponde a un control, engadir ou eliminar filas e modificar a data
    da seguinte visita.

    \item Se elimina unha fila por erro, pode desfacer a eliminación.

    \item A persoa usuaria finaliza as correccións.

    \item A aplicación comproba que exista polo menos unha fila, que as datas
    sexan válidas e non estean repetidas, que os días que non sexan de control
    conteñan unha dose válida e que a seguinte visita teña unha data válida.

    \item A persoa usuaria confirma definitivamente a pauta.

    \item Se o actor é o paciente, a nova información substitúe a pauta local. Se se trata da mesma folla
    xa existente, consérvanse os estados das tomas previamente rexistrados.

    \item A aplicación reprograma os recordatorios aplicables e informa aos
    supervisores vinculados da existencia da nova folla.

    \item Se o actor é un supervisor, a pauta revisada non se garda como
    tratamento propio do seu dispositivo. A información envíase cifrada ao
    paciente vinculado.

    \item O dispositivo do paciente recompón a información recibida e só a
    incorpora cando dispón de todos os datos necesarios.
\end{enumerate}

\item[Fluxos alternativos:] Se a validación detecta algún erro, non se permite
confirmar ata corrixilo. O actor pode cancelar as correccións e recuperar a
última versión revisada. Se abandona a revisión sen confirmar, a pauta anterior
permanece intacta. Se un supervisor non pode completar o envío ao paciente, a
operación non se considera finalizada e non se incorpora unha pauta incompleta.

\item[Postcondicións:] Se confirma o paciente, a pauta revisada queda gardada
localmente. Se confirma un supervisor, a pauta queda gardada no dispositivo do
paciente unha vez recibida completamente. Unha revisión cancelada non modifica
a pauta existente.

\item[Requisitos relacionados:] RF-11, RF-12, RF-13, RF-18, RF-19, RF-21 e
RF-24.

\end{description}


% ============================================================
% CU-06
% ============================================================

\section{CU-06: Consultar o tratamento actual}
\label{cu:consultar-tratamento}

\begin{description}

\item[Obxectivo:] Permitir ao paciente consultar desde a pantalla principal a
indicación correspondente ao tratamento, así como a pauta gardada e os datos
do próximo control cando corresponda.

\item[Actor principal:] Paciente.

\item[Actor secundario:] Firebase Cloud Messaging, cando se comunica aos
supervisores a existencia dunha nova toma esquecida.

\item[Disparador:] O paciente accede á pantalla principal ou volve á aplicación.

\item[Precondicións:] A aplicación está configurada co rol de paciente. A
consulta da información xa almacenada localmente non require conexión á rede.

\item[Fluxo principal:] \hfill
\begin{enumerate}

    \item Cárganse a pauta, os estados das tomas e os datos xerais da folla
    almacenados localmente.

    \item A aplicación busca a entrada correspondente ao día actual. Se non
    existe, selecciona a primeira entrada futura dispoñible.

    \item Se a entrada corresponde a unha toma, móstranse a dose, a data, a hora
    habitual e o seu estado.

    \item Se a entrada corresponde a un día de control, móstrase a información
    do control no lugar dunha toma.

    \item No modo normal móstranse ademais a pauta completa e os datos do próximo
    control ou visita.
\end{enumerate}

\item[Fluxos alternativos:] Se non existe ningunha pauta gardada, móstrase unha
mensaxe indicando que é necesario incorporar unha folla de tratamento e
ofrécese acceso ao fluxo de dixitalización. Se non existe unha entrada para hoxe
pero si existen entradas futuras, móstrase a primeira delas. Se non quedan
entradas actuais nin futuras, indícase que non hai tomas pendentes. Se está
activado o modo sinxelo, móstrase unicamente a información esencial da toma ou
do control e ocúltanse a pauta, a información complementaria e a navegación
inferior.

\item[Postcondicións:] Ningunha.

\item[Requisitos relacionados:] RF-14 e RF-19.

\end{description}


% ============================================================
% CU-07
% ============================================================

\section{CU-07: Confirmar a toma do día}
\label{cu:confirmar-toma}

\begin{description}

\item[Obxectivo:] Rexistrar que o paciente realizou a toma correspondente ao día
actual e conservar información sobre o momento no que foi confirmada.

\item[Actor principal:] Paciente.

\item[Actor secundario:] Firebase Cloud Messaging, cando existen supervisores
vinculados.

\item[Disparador:] O paciente preme a opción de confirmar a toma.

\item[Precondicións:] Existe unha entrada correspondente ao día actual, contén
unha dose que debe tomarse, non é un día de control e a toma continúa pendente.

\item[Fluxo principal:] \hfill
\begin{enumerate}
    \item A aplicación mostra un diálogo coa dose que se vai rexistrar e solicita
    confirmación.

    \item O paciente confirma a operación.

    \item A aplicación obtén a hora actual do dispositivo.

    \item Calcúlase a diferenza en minutos respecto da hora habitual configurada.

    \item Se a confirmación se realiza despois da hora habitual, a toma queda
    identificada como realizada fóra de hora. En caso contrario queda
    identificada como realizada.

    \item Gárdanse localmente o estado da toma, o instante exacto da confirmación
    e a desviación respecto da hora habitual.

    \item A pantalla principal actualízase para mostrar o novo estado.

    \item Se existen supervisores vinculados, envíaselles a actualización da
    toma.

    \item Cancélase o recordatorio posterior asociado a unha posible toma non
    confirmada durante esa xornada.
\end{enumerate}

\item[Fluxos alternativos:] Se o paciente cancela o diálogo, non se realiza
ningún cambio. A opción de confirmación non está dispoñible para días futuros,
días de control, días nos que non corresponde tomar medicación ou tomas que xa
non estean pendentes. Un fallo na comunicación cos supervisores non desfai a
confirmación gardada localmente.

\item[Postcondicións:] A toma queda rexistrada como realizada ou realizada fóra
de hora, xunto coa hora exacta e a desviación correspondente.

\item[Requisitos relacionados:] RF-15, RF-18 e RF-19.

\end{description}


% ============================================================
% CU-08
% ============================================================

\section{CU-08: Consultar calendario e cumprimento}
\label{cu:consultar-calendario-cumprimento}

\begin{description}

\item[Obxectivo:] Permitir ao paciente visualizar a distribución temporal da
pauta e consultar un resumo do seu cumprimento.

\item[Actor principal:] Paciente.

\item[Disparador:] O paciente accede á pantalla do calendario.

\item[Precondicións:] A aplicación está configurada co rol de paciente. A
consulta dos datos gardados non require conexión.

\item[Fluxo principal:] \hfill
\begin{enumerate}

    \item Carga a pauta, os estados das tomas e a data da seguinte visita.

    \item Agrupa as entradas da pauta polos meses aos que pertencen.

    \item Mostra un calendario no que se diferencian as doses programadas, os
    días sen toma, os días de control e os diferentes estados das tomas.

    \item Móstrase a data da próxima cita.Indícase tamén cantos días faltan,
    se a cita é hoxe ou se xa pasou.

    \item A aplicación calcula o número de tomas realizadas.

    \item Calcula o número de tomas correspondentes a días anteriores que non
    foron confirmadas.

    \item A partir destes valores calcula a porcentaxe de cumprimento e mostra
    un resumo xunto cun indicador gráfico.
\end{enumerate}

\item[Fluxos alternativos:] Se non existe ningunha pauta gardada, indícase que
non hai datos dispoñibles. Se non existe unha seguinte visita ou a súa data non
pode ser interpretada, non se mostra o número de días restante. Se aínda non
existen tomas realizadas nin esquecidas, a porcentaxe de cumprimento é cero.

\item[Postcondicións:] Ningunha

\item[Requisitos relacionados:] RF-16.

\end{description}


% ============================================================
% CU-09
% ============================================================

\section{CU-09: Consultar evolución do tratamento}
\label{cu:consultar-evolucion}

\begin{description}

\item[Obxectivo:] Permitir ao paciente consultar a evolución dos principais
valores asociados ao tratamento e información sobre a puntualidade das tomas.

\item[Actor principal:] Paciente.

\item[Disparador:] O paciente accede á pantalla de progreso.

\item[Precondicións:] A aplicación está configurada co rol de paciente. A
consulta utiliza os datos almacenados localmente e non require conexión.

\item[Fluxo principal:] \hfill
\begin{enumerate}
    \item A aplicación carga o histórico asociado ás follas de tratamento.

    \item Recupera o INR actual e a dose semanal actual.

    \item Recupera os rexistros das tomas que permiten analizar o seu
    cumprimento horario.

    \item Mostra o INR e a dose semanal actuais.

    \item Representa graficamente a evolución do INR utilizando os datos
    históricos e o valor actual.

    \item Representa graficamente a evolución da dose semanal.

    \item Calcula a desviación media respecto da hora habitual das tomas para as
    que existe esta información.

    \item Conta o número de tomas confirmadas despois da hora habitual.

    \item Mostra os rexistros recentes de confirmación xunto coa hora e a
    desviación correspondente.
\end{enumerate}

\item[Fluxos alternativos:] Se non existen valores suficientes para construír
algunha das gráficas, indícase que non hai datos suficientes. Se non existen
rexistros con información horaria, a desviación media non se mostra.

\item[Postcondicións:] A consulta non modifica a pauta nin os estados das tomas.

\item[Requisitos relacionados:] RF-17.

\end{description}


% ============================================================
% CU-10
% ============================================================

\section{CU-10: Supervisar os pacientes vinculados}
\label{cu:supervisar-pacientes}

\begin{description}

\item[Obxectivo:] Permitir ao supervisor consultar a relación dos seus pacientes
vinculados e obter unha visión resumida do estado de cada un.

\item[Actor principal:] Supervisor.

\item[Actor secundario:] Firebase Cloud Messaging.

\item[Disparador:] O supervisor accede á súa pantalla principal.

\item[Precondicións:] A aplicación está configurada co rol de supervisor. Para
actualizar os datos dos pacientes é necesaria conexión cos dispositivos
vinculados.

\item[Fluxo principal:] \hfill
\begin{enumerate}
    \item A aplicación carga os pacientes vinculados almacenados localmente.

    \item Solicita aos dispositivos dos pacientes unha actualización do seu
    estado.

    \item Os pacientes que reciben a solicitude responden coa información
    actualizada do tratamento.

    \item A aplicación actualiza os rexistros locais a medida que recibe as
    respostas.

    \item Para cada paciente móstrase o seu nome ou, se non está dispoñible, un
    identificador xenérico.

    \item Se existe unha folla de tratamento, móstranse a dose e o estado do día
    actual, o número de días tomados e non tomados, a porcentaxe de cumprimento
    e a seguinte cita.

    \item O supervisor pode solicitar manualmente unha nova actualización.

    \item A aplicación realiza tamén novas solicitudes de sincronización ao
    volver ao primeiro plano e periodicamente mentres permanece activa.

    \item Ao seleccionar un paciente cos datos básicos configurados, accédese
    ao CU-11.

    \item Se ao paciente lle falta o nome ou a hora habitual, accédese previamente
    ao CU-12.
\end{enumerate}

\item[Fluxos alternativos:] Se non existen pacientes vinculados, indícase esta
situación e ofrécese a posibilidade de realizar unha nova vinculación mediante
o CU-02. Se un paciente aínda non dispón dunha folla, móstrase unha indicación
para incorporala. Se falla a sincronización, consérvase e móstrase a última
información dispoñible localmente.

\item[Postcondicións:] O supervisor conserva a información máis recente que
recibiu de cada paciente. A consulta non modifica a pauta dos pacientes.

\item[Requisitos relacionados:] RF-05, RF-19 e RF-20.

\end{description}


% ============================================================
% CU-11
% ============================================================

\section{CU-11: Consultar o detalle dun paciente}
\label{cu:consultar-detalle-paciente}

\begin{description}

\item[Obxectivo:] Permitir ao supervisor consultar información detallada do
tratamento dun paciente vinculado.

\item[Actor principal:] Supervisor.

\item[Actor secundario:] Firebase Cloud Messaging, cando se recibe unha
actualización do paciente durante a consulta.

\item[Disparador:] O supervisor selecciona un paciente na súa lista.

\item[Precondicións:] O paciente está vinculado ao supervisor e dispón dos datos
básicos necesarios para acceder á súa vista de detalle.

\item[Fluxo principal:] \hfill
\begin{enumerate}
    \item A aplicación abre a información local correspondente ao paciente
    seleccionado.

    \item Mostra a dose semanal actual e o INR actual.

    \item Mostra a dose correspondente ao día actual, o seu estado e a hora
    habitual configurada.

    \item Representa graficamente a evolución do INR.

    \item Representa graficamente a evolución da dose semanal.

    \item Mostra a seguinte visita dispoñible.

    \item Se durante a consulta se recibe unha actualización do paciente, a
    información local é recargada.

    \item Desde esta pantalla o supervisor pode iniciar o CU-04 para incorporar
    unha nova folla de tratamento no nome do paciente.
\end{enumerate}

\item[Fluxos alternativos:] Se non existen datos suficientes para unha gráfica,
indícase que non hai información suficiente. Os valores non dispoñibles
móstranse sen dato asociado. Se non existe conexión, pódese consultar o último
estado almacenado no dispositivo. Se faltan o nome ou a hora habitual do
paciente, desde a lista accédese ao CU-12 antes de entrar nesta vista.

\item[Postcondicións:] Non se modifica o tratamento do paciente. Unicamente se
actualiza a información mostrada se se reciben novos datos sincronizados.

\item[Requisitos relacionados:] RF-21.

\end{description}


% ============================================================
% CU-12
% ============================================================

\section{CU-12: Configurar datos dun paciente supervisado}
\label{cu:configurar-paciente-supervisado}

\begin{description}

\item[Obxectivo:] Permitir ao supervisor establecer o nome e a hora habitual da
toma dun paciente vinculado cando estes datos sexan necesarios para o
seguimento.

\item[Actor principal:] Supervisor.

\item[Actor secundario:] Firebase Cloud Messaging.

\item[Disparador:] O supervisor selecciona desde a lista un paciente ao que lle
falta o nome ou a hora habitual da toma.

\item[Precondicións:] Existe unha vinculación válida co paciente e a relación
dispón da información necesaria para realizar unha comunicación segura.

\item[Fluxo principal:] \hfill
\begin{enumerate}
    \item A aplicación mostra unha pantalla coa configuración do paciente.

    \item O supervisor pode modificar o nome mostrado.

    \item O supervisor selecciona a hora habitual da toma. Inicialmente existe
    unha hora predefinida que pode ser modificada.

    \item O supervisor preme a opción de gardar e enviar.

    \item A aplicación cifra os datos empregando a información asociada á
    vinculación.

    \item A configuración envíase ao dispositivo do paciente.

    \item O paciente recibe e valida a mensaxe.

    \item O dispositivo do paciente actualiza o nome e a hora recibidos.

    \item Os recordatorios locais do paciente son programados novamente segundo
    a nova hora.

    \item O paciente envía aos seus supervisores un novo estado actualizado.

    \item O supervisor actualiza a información almacenada do paciente.
\end{enumerate}

\item[Fluxos alternativos:] Se a vinculación non dispón dunha relación segura ou
non se pode entregar a configuración, móstrase unha mensaxe de erro e o
supervisor permanece na pantalla para poder intentalo novamente.

\item[Postcondicións:] O dispositivo do paciente conserva os novos datos e os
supervisores reciben posteriormente o estado actualizado.

\item[Requisitos relacionados:] RF-07, RF-18, RF-19 e RF-24.

\end{description}


% ============================================================
% CU-13
% ============================================================

\section{CU-13: Consultar o centro sanitario}
\label{cu:consultar-centro}

\begin{description}

\item[Obxectivo:] Facilitar ao paciente o acceso aos datos de contacto do centro
sanitario indicado na folla de tratamento e permitir iniciar unha chamada.

\item[Actor principal:] Paciente.

\item[Disparador:] O paciente selecciona a opción para chamar ao centro sanitario
mostrado xunto á información do próximo control.

\item[Precondicións:] A folla de tratamento identifica un centro sanitario. É
necesaria conexión á rede para consultar os seus datos de contacto.

\item[Fluxo principal:] \hfill
\begin{enumerate}
    \item A aplicación obtén o nome do centro sanitario almacenado coa folla de
    tratamento.

    \item O paciente selecciona a opción de chamar ao centro.

    \item A aplicación solicita ao servidor a información correspondente ao
    centro indicado.

    \item O servidor busca o centro e devolve os datos dispoñibles.

    \item A aplicación obtén o número de teléfono.

    \item Se o número está dispoñible, abre o marcador telefónico do dispositivo
    con ese número preparado para realizar a chamada.
\end{enumerate}

\item[Fluxos alternativos:] Se non existe un centro identificado, non se mostra
a opción de chamada. Se o centro non pode ser atopado, non dispón dun teléfono,
se produce un erro de conexión ou non se pode abrir o marcador, móstrase unha
mensaxe de erro.

\item[Postcondicións:] A información do tratamento non sofre modificacións. Se
a operación se completa, o marcador telefónico queda aberto co número do centro.

\item[Requisitos relacionados:] RF-22.

\end{description}


% ============================================================
% CU-14
% ============================================================

\section{CU-14: Revogar unha vinculación}
\label{cu:revogar-vinculacion}

\begin{description}

\item[Obxectivo:] Permitir que calquera das dúas partes finalice unha relación
de supervisión existente.

\item[Actor principal:] Paciente ou supervisor.

\item[Actor secundario:] Firebase Cloud Messaging.

\item[Disparador:] O paciente decide quitar unha persoa supervisora desde os seus
axustes ou o supervisor decide deixar de supervisar un paciente.

\item[Precondicións:] Existe unha vinculación entre os dous dispositivos. Para
comunicar correctamente a revogación á outra parte é necesaria conexión á rede.

\item[Fluxo principal:] \hfill
\begin{enumerate}
    \item A aplicación mostra as vinculacións existentes do actor.

    \item A persoa usuaria selecciona a relación que quere eliminar.

    \item A aplicación mostra un diálogo explicando as consecuencias da
    desvinculación e solicita confirmación.

    \item A persoa usuaria confirma a operación.

    \item Se a revogación a inicia o paciente, a aplicación envía ao supervisor
    unha mensaxe segura indicando que a relación debe ser eliminada.

    \item Unha vez enviado o aviso, o dispositivo do paciente elimina localmente
    a información correspondente a ese supervisor.

    \item O supervisor que recibe a mensaxe elimina o paciente dos seus
    rexistros locais e cancela os recordatorios asociados a el.

    \item Se a revogación a inicia o supervisor, envíase de forma equivalente
    unha mensaxe segura ao paciente.

    \item Unha vez enviado o aviso, o supervisor elimina da súa lista o paciente
    e cancela os recordatorios asociados.

    \item O paciente que recibe a mensaxe elimina a información correspondente
    ao supervisor que iniciou a desvinculación.
\end{enumerate}

\item[Fluxos alternativos:] Se a persoa usuaria cancela o diálogo, a relación
permanece intacta. Se non existe información criptográfica válida para a
vinculación ou non se pode comunicar a revogación á outra parte, móstrase un
erro e a aplicación non completa a eliminación local iniciada polo usuario.

\item[Postcondicións:] A relación deixa de estar dispoñible nos dispositivos e
non pode empregarse para continuar intercambiando información. As restantes
vinculacións do paciente ou do supervisor non se ven afectadas.

\item[Requisitos relacionados:] RF-19, RF-24 e RF-25.

\end{description}